\documentclass[11pt]{article}
\usepackage[utf8]{inputenc}
\usepackage{amsmath}
\usepackage{amsfonts}
\usepackage{amssymb}
\usepackage{bm}
\usepackage{geometry}
\usepackage{xcolor}
\usepackage{booktabs}

\geometry{margin=0.8in}

\title{PRISM-EM for Non-Stationary Poisson GLMs with States from Spike Trains:\\
Model, Objective, and Implementation. Ver3c}
\author{Jan Balewski, NERSC}
\date{\today}

\begin{document}

\maketitle

\begin{abstract}
This document describes the current \texttt{prism\_EM\_train3c.py} implementation
for fitting a non-stationary
Poisson GLM with latent state mixing (PRISM-EM)~\footnote{PRISM: PRobabilistic Inference of Sparse Mixed-state dynamics; ~~EM: Expectations Maximization}. Input spike trains are read from \texttt{spikesData/<dataName>.spikes.npz} (non-stationary output of \texttt{gen\_nonStationarySpikes3c.py}). The model uses a shared connectivity matrix across states, state-specific bias vectors, and time-varying simplex coefficients. Training is performed by block coordinate descent: a sequential projected-gradient E-step for coefficients followed by Viterbi decoding to one-hot state assignments, and an Adam-based M-step for model parameters with off-diagonal L1 shrinkage and spectral-radius projection. Spectral-radius enforcement, L1 pruning, and learning-rate decay are each activated after a configurable number of EM iterations. Distributed training (\texttt{torchrun}) shards the E-step over time bins and synchronizes \(A\), \(B\), and \(c\) across GPUs. We summarize the exact objective used in code, the decoding method, all current default hyperparameters, initialization modes (data-driven vs random) for \(A\), \(B\), and \(C\), and the saved fit file layout.
\end{abstract}

\tableofcontents
\newpage

% ---------------------------------------------------------------
\section{EM Training for Non-Stationary Poisson GLM with states}
\label{sec:em}
% ---------------------------------------------------------------

\subsection{Introduction and Model Assumptions}

We consider a population of \(N\) neurons observed over \(T\) discrete time
bins of width \(\Delta t\) (seconds). The population activity is governed
by \(M\) latent dynamical states that switch over time. The key
structural assumption is that all states share a single
connectivity matrix \(A \in \mathbb{R}^{N \times N}\) but differ in
baseline excitability through per-state bias vectors
\(B_m \in \mathbb{R}^N,\; m = 0, \dots, M-1\) (0-based state indices in code).

At each time bin \(t\), a coefficient vector
\(c_t \in \mathbb{R}^M\), constrained to the probability simplex
(\(c_{mt}\ge 0,\; \sum_m c_{mt}=1\)), selects the effective bias as a
convex combination:
\begin{equation}
  B_t = \sum_{m=1}^{M} c_{mt}\, B_m = c_t^\top \mathbf{B},
\end{equation}
where \(\mathbf{B} \in \mathbb{R}^{M \times N}\) stacks the \(M\) bias
vectors row-wise.

Spike counts follow
\begin{equation}
  \eta_t = A\,Y_{t-1} + B_t,\qquad
  \tilde{\eta}_t = \min(\eta_t,\eta_{\mathrm{clip}}),\qquad
  \mu_t = \exp(\tilde{\eta}_t)\,\Delta t,\qquad
  Y_t \sim \mathrm{Poisson}(\mu_t).
  \label{eq:glm}
\end{equation}

\paragraph{Notation summary.}
\(N\): number of neurons; 
\(M\): number of latent states; 
\(T\): number of time bins;
\(\Delta t\): time-bin width (from metadata);
\(\eta_{\mathrm{clip}}\): clipping threshold (from metadata);
\(Y_t \in \mathbb{N}_0^N\): observed spikes;
\(\mu_t \in \mathbb{R}_+^N\): expected spike count per bin.

% ---------------------------------------------------------------
\subsection{Objective Function}

The training objective uses uniformly weighted (unweighted) Poisson NLL in both the E-step and M-step:
\begin{equation}
\mathcal{J}
=
\underbrace{\sum_{t=2}^{T}\sum_{i=1}^{N}
\left[\mu_{t,i}-Y_{t,i}\left(\tilde{\eta}_{t,i}+\log\Delta t\right)\right]}_{\text{Poisson NLL}}
+
\underbrace{\lambda_2\sum_{t=2}^{T}\|\vec c_t- \vec c_{t-1}\|^2}_{\text{temporal smoothness (E-step only)}}
+
\underbrace{\lambda_3\,\overline{|A_{\mathrm{off}}|}}_{\text{L1 sparsity on off-diagonal (M-step only)}},
\label{eq:objective}
\end{equation}
with simplex constraints \(c_t\in\Delta^{M-1}\) and spectral constraint
\(\rho(A)\le \rho_{\max}\), where the off-diagonal L1 term is
\begin{equation}
  \overline{|A_{\mathrm{off}}|}
  =
  \frac{1}{N(N-1)}
  \sum_{\substack{i,j=1\\i\neq j}}^{N}|A_{ij}|.
  \label{eq:l1_offdiag}
\end{equation}
In the $\lambda_2$ term, the sum over states $m$
is already implicit in the $L_2$ norm.


\subsection{Why Block Coordinate Descent (EM) over Joint Optimization}

In principle one could jointly optimize \(A\), \(\mathbf{B}\), and
\(c_{1:T}\) by simultaneous gradient descent on the full objective
\(\mathcal{J}\).  In practice, the EM-style alternation is preferable
for three structural reasons.

\paragraph{Breaking the bilinear coupling.}
The coefficients \(c_t\) and the bias matrix \(\mathbf{B}\) enter the
log-rate linearly through the product \(c_t^\top\mathbf{B}\), which
then passes through an exponential.  When both are free to move, the
landscape contains saddle points and flat valleys: a large \(c_{tm}\)
paired with a small \(B_m\) produces the same likelihood as the
reverse.  Fixing one side at each step eliminates this ambiguity.  The
Viterbi hard-assignment variant sharpens the separation further:
with one-hot \(c_t = e_{\hat{S}_t}\), each time bin selects exactly one
bias row, so the M-step reduces to \(M\) independent Poisson GLM
regressions and the bilinear degeneracy vanishes entirely.

\paragraph{Exploiting temporal structure in the E-step.}
The coefficient vectors \(c_{1:T}\) number \(O(TM)\) and vastly
outnumber the \(O(N^2+MN)\) model parameters.  Joint optimization
would need to balance learning rates across these disparate groups
while enforcing a simplex constraint at every time bin.  By contrast,
the E-step decomposes into a single forward sweep of cheap
\(M\)-dimensional simplex-projected gradient updates, naturally
incorporating the temporal smoothness penalty
\(\lambda_2\|c_t-c_{t-1}\|^2\) through the sequential dependence on the
previous accepted \(c_{t-1}\).

\paragraph{Avoiding the log-sum-exp mismatch.}
Under soft (probabilistic) state assignments the M-step must fit
parameters to the mixture rate
\(\sum_m c_{tm}\exp(\eta_{tm})\), whereas the model parameterizes the
log-rate as \(\sum_m c_{tm}\,\eta_{tm}\) inside a single exponential.
These two quantities differ whenever the assignment is not one-hot,
creating a systematic bias in the gradient signal for \(\mathbf{B}\).
The hard-EM strategy—Viterbi decode followed by one-hot
encoding—eliminates this mismatch by construction, at the cost of
losing the monotonic likelihood improvement guarantee of soft EM, a
trade-off that is empirically favorable.

% ---------------------------------------------------------------
\subsection{Algorithm: Block Coordinate Descent (EM)}

The non-convex problem is solved by alternating:
\begin{itemize}
  \item E-step: update \(c_{1:T}\) with \(A,\mathbf{B}\) fixed, then Viterbi-decode to one-hot state assignments.
  \item M-step: update \(A,\mathbf{B}\) with one-hot \(c_{1:T}\) fixed.
\end{itemize}
for \(K_{\mathrm{EM}}\) outer iterations.

\subsubsection{E-step: Coefficient Inference via Simplex PGD, then One-Hot Encoding}

Holding \((A,\mathbf{B})\) fixed, coefficients are inferred by one forward
sweep from \(t=2\) to \(T\). At each time bin:
\begin{enumerate}
  \item Build \(Z_t[m,:] = A\,Y_{t-1}+B_m\in\mathbb{R}^{M\times N}\).
  \item Run \texttt{pgd\_iter} projected-gradient iterations:
  \[
    \eta_t = c_t^\top Z_t,\quad
    \tilde{\eta}_t = \min(\eta_t,\eta_{\mathrm{clip}}),\quad
    \mu_t = \exp(\tilde{\eta}_t)\Delta t,
  \]
  \[
    g_t = Z_t\left(\mu_t-Y_t\right) + 2\lambda_2(c_t-c_{t-1}),
    \qquad
    c_t \leftarrow \Pi_{\Delta^{M-1}}(c_t-\alpha_E g_t).
  \]
  \item Accept \(c_t\), continue to next time bin.
\end{enumerate}
Projection uses the sort-based simplex projection (Duchi et al., \(O(M\log M)\)).
Bin \(t=0\) is not updated; the sweep starts at \(t=1\) (first lag pair
\((Y_0,Y_1)\)), matching the objective sum over \(t=2,\dots,T\) in
1-based notation.
In the distributed setting, each rank updates a disjoint time shard of
\(c_{1:T}\); shard results are averaged where multiple ranks overlap at
shard boundaries, then broadcast so all ranks share the same \(c\).

After the forward sweep completes, the soft coefficients \(c_{1:T}\) are decoded into a hard state sequence via Viterbi (Section~\ref{sec:decode}). The M-step then receives one-hot encoded state assignments rather than the soft simplex probabilities:
\begin{equation}
  \hat{S}_{1:T} = \mathrm{Viterbi}(c_{1:T}), \qquad
  c_t^{\mathrm{(M\text{-}step)}} = e_{\hat{S}_t},
\end{equation}
where \(e_m\) is the \(m\)-th standard basis vector. This ensures the M-step fits each state's parameters from cleanly separated data rather than from soft mixtures.

\subsubsection{M-step: Dictionary Update via Adam}

Holding \(c_t\) fixed (as one-hot vectors), update \(A,\mathbf{B}\) for \(K_M\) epochs using
Adam over shuffled mini-batches of \((Y_{t-1},Y_t,c_t)\).

Per mini-batch:
\begin{equation}
  \ell_{\mathrm{M}} =
  \frac{1}{|\mathcal{B}|}
  \sum_{t\in\mathcal{B}}\sum_{i=1}^{N}
  \left[-Y_{t,i}\log(\mu_{t,i}+\varepsilon)+\mu_{t,i}\right]
  + \lambda_3\,\overline{|A_{\mathrm{off}}|},
\end{equation}
with \(\varepsilon=10^{-8}\).

Then, subject to delay scheduling (operations are activated only after a configurable number of EM iterations have passed):
\begin{enumerate}
  \item Off-diagonal soft-thresholding (activated after \(K_{\mathrm{prune}}\) EM iterations):
  \[
    A_{ij}\leftarrow \operatorname{sign}(A_{ij})
    \left(|A_{ij}|-\alpha_M\lambda_3\right)_+,\quad i\neq j.
  \]
  \item Periodic spectral-radius projection (activated after \(K_{\rho}\) EM iterations, applied every \(R\) batches). On each application, in the multi-GPU setting, \(A\) and \(B\) are first averaged across ranks, then the spectral constraint is enforced, and the result is broadcast:
  \[
    A \leftarrow A\cdot \min\!\left(1,\frac{\rho_{\max}}{\rho(A)}\right).
  \]
\end{enumerate}

Learning rate follows linear decay from \(\alpha_M^{(0)}\) to
\(\alpha_M^{(0)}f_{\mathrm{end}}\). The decay is activated only after \(K_{\mathrm{lr}}\) EM iterations have passed; the total number of decay steps spans the remaining M-step epochs.

% ---------------------------------------------------------------
\subsection{State Decoding}
\label{sec:decode}

Viterbi decoding is applied both within the EM loop (after each E-step, to produce one-hot assignments for the M-step) and after the final EM iteration (to produce the output state sequence). The decoder uses a homogeneous transition prior with
\(\tau_{\mathrm{dwell}}=\texttt{decode\_dwell\_sec}\):
\begin{equation}
  p_{\mathrm{stay}}=\exp\!\left(-\frac{\Delta t}{\tau_{\mathrm{dwell}}}\right),
  \qquad
  p_{\mathrm{switch}}=\frac{1-p_{\mathrm{stay}}}{M-1}.
\end{equation}
The decoded sequence is
\begin{equation}
  \hat{S}_{1:T}=\arg\max_{\{s_t\}}
  \sum_t\left[\log c_{t,s_t}+\log p(s_t\mid s_{t-1})\right].
\end{equation}
Per-bin confidence is
\[
\mathrm{CL}_t = 1-\max_m c_{t,m}.
\]

% ---------------------------------------------------------------
\section{Input and Output File Pattern}
\label{sec:io}
% ---------------------------------------------------------------

The ordinary EM trainer is \texttt{prism\_EM\_train3c.py}. It is self-contained
through \texttt{PrismEM\_Workhorse3c.py}; no code from older
\texttt{ver3} trainers is required.

\paragraph{Input.}
The required input spike file is
\[
\texttt{<basePath>/spikesData/<dataName>.spikes.npz}.
\]
The file must contain \texttt{spikes} with shape \((T,N)\) and
\texttt{single\_rates} with shape \((N,)\). Metadata must provide
\texttt{time\_step\_sec} and \texttt{poisson\_eta\_clip}. Synthetic files
from \texttt{gen\_nonStationarySpikes3c.py} also carry provenance metadata,
which is copied into the fit output.

For real experimental recordings, the preprocessing program is
\texttt{prep\_bioexp3c.py}. Its output arrays are compatible with the EM
trainer, but the output directory must be chosen to match the EM file
layout. Run preprocessing with
\[
\texttt{--dataPath <basePath>/spikesData}
\]
and choose \texttt{--shortName <dataName>} if an explicit EM data name is
desired. The resulting files are
\[
\texttt{<basePath>/spikesData/<dataName>.spikes.npz},
\qquad
\texttt{<basePath>/spikesData/<dataName>.bioExp.npz}.
\]
The \texttt{.spikes.npz} file is the direct EM input. The companion
\texttt{.bioExp.npz} file is not read during training, but it is expected by
the biological-data evaluation and plotting path to sit next to the spike
file with the same stem. In the spike metadata,
\texttt{provenance.experiment\_name} must therefore match
\texttt{<dataName>}.

\paragraph{Output.}
Rank~0 creates \texttt{<basePath>/prismFit/} if needed and writes
\[
\texttt{<basePath>/prismFit/<fitName>.prismEM.npz}.
\]
If \texttt{--fitName} is omitted, the program generates
\[
\texttt{<dataName>\_<hash4>}.
\]
The metadata provenance field \texttt{EMtrain\_file} is set to the same
output stem.

\paragraph{Canonical commands.}
\begin{verbatim}
torchrun --standalone --nnodes=1 --nproc_per_node=4 \
  ./prism_EM_train3c.py \
  --basePath $basePath \
  --dataName <dataName> \
  --num_states <M> \
  --time_range_sec 0 80

./prism_EM_train3c.py \
  --basePath $basePath \
  --dataName <dataName> \
  --num_states <M> \
  --fitName <fitName>
\end{verbatim}

The first form uses distributed GPU training when launched under
\texttt{torchrun}; the second form runs on a single process, using CUDA if a
GPU is visible and CPU otherwise.

% ---------------------------------------------------------------
\section{Default Hyperparameters and Runtime Settings}
\label{sec:defaults}
% ---------------------------------------------------------------

\begin{table}[h]
\centering
\begin{tabular}{lll}
\toprule
Group & Parameter (\texttt{CLI}) & Default \\
\midrule
EM structure & \texttt{--num\_em\_iters} & 20 \\
 & \texttt{--m\_epochs} & 2 \\
E-step & \texttt{--pgd\_iter} & 5 \\
 & \texttt{--lr\_estep} & 0.03 \\
 & \texttt{--lambda2} & 2.0 \\
M-step & \texttt{--lr\_mstep} & \(3 \times 10^{-3}\) \\
 & \texttt{--lr\_end\_factor} & 0.1 \\
 & \texttt{--lambda3} & 0.02 \\
 & \texttt{--rho\_max} & 0.95 \\
 & \texttt{--prescale\_m\_step\_4\_ArhoMax} & 120 (batches) \\
 & \texttt{--delay\_em\_iter\_4\_ArhoMax} & [3] \\
 & \texttt{--delay\_em\_iter\_4\_lrDecay} & [1] \\
 & \texttt{--delay\_em\_iter\_4\_Aprune} & 1 \\
 & \texttt{--batch\_size} & 2048 \\
Data window & \texttt{--time\_range\_sec} & [0.0, 60.0] s \\
Decoding & \texttt{--decode\_dwell\_sec} & 0.3 s \\
Init & \texttt{--init\_A} & data (or rand)\\
 & \texttt{--init\_B} & data (or rand) \\
 & \texttt{--init\_states} & data (or rand) \\
\bottomrule
\end{tabular}
\caption{Current defaults in \texttt{PrismEM\_Workhorse3c.py}, as used by
\texttt{prism\_EM\_train3c.py}. The two delay arguments may be given as one
or two integers: start iteration, and optionally target iteration. If the
target is omitted, the code uses \texttt{num\_em\_iters}.}
\end{table}

\paragraph{Required input.}
\texttt{--num\_states} (\texttt{-M}) is required (no default).
\texttt{--dataName} and \texttt{--basePath} locate input spikes under
\texttt{<basePath>/spikesData/}; fits are written to
\texttt{<basePath>/prismFit/<fitName>.prismEM.npz}.

\paragraph{Dataset-provided values.}
\(\Delta t\) (\texttt{time\_step\_sec}) and \(\eta_{\mathrm{clip}}\)
(\texttt{poisson\_eta\_clip}) are loaded from spike-file metadata.

% ---------------------------------------------------------------
\section{Initialization of Model}
\label{sec:init}
% ---------------------------------------------------------------

Initialization is controlled by:
\[
\texttt{--init\_states},\quad \texttt{--init\_B},\quad  \texttt{--init\_A}
\in \{\texttt{data},\texttt{rand}\},
\]
with default \texttt{data} for all three.

\subsection{State-Coefficient Initialization (\(C\equiv\{c_t\}\))}

\paragraph{random mode:}
\[
c_t = \frac{1}{M}\mathbf{1},\quad S_t=-1 \;\; \forall t.
\]
So coefficients start uniform over states.

\paragraph{data mode:}
\begin{enumerate}
  \item Aggregate spikes into coarse blocks of length
  \(\lceil\tau_{\mathrm{dwell}}/\Delta t\rceil\) bins, where
  \(\tau_{\mathrm{dwell}}=\texttt{decode\_dwell\_sec}\) (default 0.3\,s).
  \item For each block, compute a scalar mean population firing rate.
  \item Find \(M-1\) thresholds minimizing within-cluster variance
  (dynamic-programming segmentation on unique rate values).
  \item Assign each block to a state by thresholds.
  \item Convert block states to soft probabilities:
  dominant state gets \(0.8\), remainder \(0.2\) split across others,
  with transition smoothing (\(2/3\) vs \(1/3\) at boundaries).
  \item Map block-level \(c\) and \(S\) back to original bin resolution.
\end{enumerate}

\subsection{Bias Initialization (\(B\))}

\paragraph{random mode:}
No external seed is applied; model keeps random Gaussian initialization from
\texttt{torch.randn(...)*0.1}.

\paragraph{data mode:}
For each neuron \(n\):
\[
r_n = \frac{1}{T}\sum_t Y_{t,n}\,/\,\Delta t,\qquad
b_n = \log(\max(r_n,10^{-6})).
\]
The same \(b\)-vector is assigned to every state row of \(B\).

\subsection{Connectivity Initialization (\(A\))}

\paragraph{random mode:}
No external seed is applied; model keeps random Gaussian initialization from
\texttt{torch.randn(...)*0.1}.

\paragraph{data mode (covariance/OLS):}
Using first \(T_{\max}=50{,}000\) bins (or fewer if shorter):
\[
A_{\mathrm{ols}}
=
\left(\sum_t Y_tY_{t-1}^\top\right)
\left(\sum_t Y_{t-1}Y_{t-1}^\top\right)^{\dagger}.
\]
Then all elements of \(A_{\mathrm{ols}}\) are scaled by a factor of 3:
\[
A \leftarrow 3\,A_{\mathrm{ols}}.
\]
No spectral-radius rescaling is applied at initialization; the spectral constraint is enforced later during training via the delayed projection mechanism.

Diagnostics (condition number, spectral radius, one-step \(R^2\), Frobenius norm,
bins used) are computed on the pre-scaled OLS estimate and saved in metadata
(\texttt{init\_A} sub-dictionary of the output metadata).

% ---------------------------------------------------------------
\section{Output File and Evaluation}
\label{sec:output}
% ---------------------------------------------------------------

After training, rank~0 writes \texttt{<basePath>/prismFit/<fitName>.prismEM.npz}.
If \texttt{--fitName} is omitted, the name is
\texttt{<dataName>\_<hash4>}.
Key arrays (shapes for \(T\) time bins, \(N\) neurons, \(M\) states):
\begin{itemize}
  \item \texttt{A\_init}, \texttt{A\_hat}: \((N,N)\);
  \item \texttt{B\_init}: \((N,)\) --- first row of the initialized \(B\) matrix
        (common across states in \texttt{data} mode);
  \item \texttt{B\_hat}: \((M,N)\) --- full per-state bias matrix;
  \item \texttt{c\_init}, \texttt{c\_hat}: \((T,M)\);
  \item \texttt{S\_init}, \texttt{S\_hat}: \((T,)\) integer state labels
        (\texttt{S\_init} is \(-1\) in \texttt{rand} init mode);
  \item \texttt{S\_hat\_CL}: \((T,)\) per-bin confidence
        \(\mathrm{CL}_t = 1-\max_m c_{t,m}\);
  \item \texttt{e\_nll\_em}, \texttt{m\_loss\_epoch}, \texttt{m\_nll\_epoch},
        \texttt{m\_l1\_epoch}, \texttt{rho\_epoch},
        \texttt{rho\_correction\_strength\_epoch},
        \texttt{nz\_edges\_epoch}, \texttt{learning\_rates}: training history;
  \item \texttt{freq\_h1d}: coarse firing-rate trace used by
        \texttt{init\_states=data};
  \item \texttt{single\_rates}: \((N,)\), copied from the input spike file;
  \item \texttt{neuron\_type}, \texttt{neuron\_Sedge}, \texttt{A\_prune}:
        source-column Dale-style postprocessing derived from fitted
        \texttt{A\_hat}. Here \texttt{neuron\_Sedge} is the off-diagonal
        column sum for each source neuron, \texttt{neuron\_type} is inferred
        from that source-column sum, and \texttt{A\_prune} preserves
        positive outgoing columns for excitatory sources and negative
        outgoing columns for inhibitory sources.
\end{itemize}
Metadata records all CLI hyperparameters, time-window bins, initialization
dictionaries, and \texttt{mstep\_state\_mode=viterbi\_onehot\_prevbin}.
Suggested evaluation:
\begin{verbatim}
./prism_EM_eval3c.py --basePath $basePath --dataName <fitName> -p a e f g
\end{verbatim}

\end{document}
